\cotefile{136}
\scanpage{282}
\mspage{- 35 -}\hfill IX
% archive sheet number: 282
% mounting strip and printed label: 353

On a un isomorphisme canonique (compatible aux opérations de $\handwritten{\pi}$) :
\[
H^{1}(A_{\overline{K}},\mathbb{Z}_{\handwritten{\ell}})\simeq\operatorname{Hom}(T_{\handwritten{\ell}}(A(\overline{K})),\mathbb{Z}_{\handwritten{\ell}})\ ,\tag*{(*)}
\]
qui montre que l'assertion 3.7 équivaut à la même assertion pour les opérations de $I$ sur le premier membre de (*). Or cette assertion a été prouvée dans III \qquad (en utilisant la résolution des singularités des schémas excellents de dimension 2).

\uline{Corollaire} 3.8. \uline{Les conditions équivalentes} (i) \uline{à} (iv) \uline{de} 3.5 \uline{équivalent aussi à la condition suivante} :

(vi) \uline{L'action de} $I$ \uline{sur} $U=T_{\handwritten{\ell}}(A(\overline{K}))$ \uline{est unipotente}.

En effet, pour prouver que (v) implique (vi), on note que si un groupe $I$ opérant de façon unipotente sur un vectoriel $E$ de dimension finie sur un corps de caractéristique nulle (en l'occurrence $\mathbb{Q}_{\handwritten{\ell}}$) admet un sous-groupe d'indice fini $I'$ opérant avec l'échelon $m$ (en l'occurrence $m=2$), alors $I$ lui-même opère avec l'échelon $m$. En effet, $I'$ étant d'indice fini dans $I$, le groupe algébrique $G'\handwritten{\subset}\text{\uline{Gl}}(E)$ qu'il définit est d'indice fini dans celui, $G$, défini par $I$, or $G$ étant unipotent est connexe (caractéristique nulle !), donc $G=G'$, donc $I$ et $I'$ ont même ``échelon d'unipotence''.

On obtient également une réciproque à 3.3 :

\uline{Corollaire} 3.9. \uline{Les notations sont celles de} 3.2. \uline{Les conditions} (i) \uline{et} (ii) \uline{de} 3.2 \uline{équivalent aussi à la condition suivante}\handwritten{$\unsure{\backslash}$} :

(iii) \uline{Pour toute extension étale} $K'$ \uline{de} $K$, \uline{désignant par} $S'$ \uline{le normalisé de} $S$ \uline{dans} $K'$, \uline{par} $A'$ \uline{le modèle de Néron de} $A_{K'}=A\otimes_{S}K'=A_{K}\otimes_{K}K'$, \uline{le morphisme canonique}
